\section{Experimental Setup}

This section describes the apparatus used to produce, cool, and trap a sample of $^{87}$Rb atoms.\CCgen{capitolo 2}{il capitolo cita pochissimo}{Contato sul sorgente: il capitolo 2 ha \textbf{13161 parole, 104 equazioni e 18 citazioni}. Il capitolo 1, con 9817 parole, ne ha 70. La sezione sul breadboard superiore, 2143 parole e 6 figure, non ne ha \textbf{nessuna}; quella sul modello numerico ne ha 3 in 42 equazioni. Per la descrizione di un apparato costruito in laboratorio e' in parte fisiologico, ma non del tutto: le tecniche usate hanno una letteratura (l'analisi con la deviazione di Allan, il mode matching su fibra, la caratterizzazione del PER, la stima del numero di atomi per fluorescenza) e citarla serve a due cose, dire da dove viene il metodo e togliere all'esaminatore il dubbio che sia stato inventato sul momento. Aggiungerne una per tecnica, non una per frase.} The atomic sample is prepared via a standard Magneto-Optical Trap (MOT) followed by an optical molasses phase, providing the cold atomic cloud required to load the Optical Dipole Trap (ODT). Spatial confinement and initial cooling rely on three pairs of counter-propagating, orthogonally polarized ($\sigma^+$/$\sigma^-$), red-detuned laser beams intersecting a quadrupole magnetic field. Following the MOT stage, the magnetic gradient is switched off to perform sub-Doppler cooling in an optical molasses (Sec. \ref{cap:trapping}).

\subsection{Setup Description}
\label{subsec:setup}

\begin{figure}[h!]
\centering
\includegraphics[width=0.5\linewidth]{Cap_2_figures/CAD_apparatus.png}
\caption{CAD rendering of the vacuum chamber and surrounding experimental setup.}
\label{fig:CAD_apparatus}
\end{figure}

The core of the apparatus is an Ultra-High Vacuum (UHV) chamber (Fig. \ref{fig:CAD_apparatus}) maintained at a baseline pressure of $\sim10^{-10}\,\mathrm{mbar}$ by an ion/getter combination pump (NEXTorr Z200). A heated ampoule of pure $^{87}$Rb supplies the atomic vapor, raising the pressure to $\sim10^{-8}\,\mathrm{mbar}$.

Optical access is provided by eight windows: two large CF100 viewports along the $z$-axis and four CF40 viewports at $45^\circ$, all AR-coated for 780 nm to optimize MOT beam transmission. Two uncoated CF40 viewports provide access for the second ODT beam and the ion pump. Due to reflections from the ion pump on the $x$-axis, imaging is performed through the $z$-axis viewports.

The MOT quadrupole magnetic field is generated by anti-Helmholtz coils ($110\,\mathrm{mm}$ inner diameter) producing a measured axial gradient of $1.33\,\mathrm{G\,cm^{-1}A^{-1}}$ (and a corresponding radial gradient of half this value with opposite sign) \cite{refRazzai}. Three mutually orthogonal pairs of Helmholtz shim coils cancel stray fields and are used to optimize the MOT-ODT spatial overlap, with calibration factors \cite{Marchesini}:

$$B_z \approx 1.4\ \mathrm{mG/mA}, \qquad B_x \approx 1.3\ \mathrm{mG/mA}, \qquad B_y \approx 1.5\ \mathrm{mG/mA}.$$

A Hollow-Core Photonic Crystal Fiber (HCPCF) is introduced from the bottom of the chamber via a custom mount, positioning its tip $\sim7\,\mathrm{mm}$ below the trap center. The optical layout is fiber-coupled using polarization-maintaining (PANDA) fibers \cite{Marchesini}. The MOT beams are delivered via three fiber collimators and retro-reflected with quarter-wave plates to maintain the correct polarization sequence.

\subsection{Laser System and Frequency Stabilization}

\begin{figure}[h!]
\centering
\includegraphics[width=0.6\linewidth]{Cap_2_figures/lines.png}
\caption{Level scheme of the $^{85}$Rb and $^{87}$Rb D2 line.}
\label{fig:lines}
\end{figure}

Cooling operates on the $^{87}$Rb $5^2S_{1/2}\rightarrow5^2P_{3/2}$ (D2) transition near $780\,\mathrm{nm}$ (Fig. \ref{fig:lines}). The cooling laser drives the closed $F=2\rightarrow F'=3$ transition, while a repumper laser on the $F=1\rightarrow F'=2$ transition recovers atoms that off-resonantly decay into the $F=1$ dark state. Both lasers are offset-locked to a common master laser~\cite{Marchesini}.

\subsubsection{Master laser}

The master reference is a $780\,\mathrm{nm}$ DFB laser diode stabilized via frequency-modulated saturation spectroscopy in a heated $^{85}$Rb vapor cell. It is locked to the $^{85}$Rb $(F=3\rightarrow F'=3)\times(F=3\rightarrow F'=4)$ crossover resonance, providing a stable absolute frequency located between the $^{87}$Rb cooling and repumping transitions~\cite{refRazzai}.
\subsubsection{Cooler and repumper lasers}

\begin{figure}[h!]
\centering
\includegraphics[width=0.8\linewidth]{Cap_2_figures/beat_locking.png}
\caption{Beat-note locking setup scheme~\cite{refRazzai}.}
\label{fig:beat_locking}
\end{figure}

The cooling and repumping light is generated by two $1560~\mathrm{nm}$ DFB diodes. Their outputs are fiber-combined, amplified to $\sim500~\mathrm{mW}$ by an Erbium-Doped Fiber Amplifier (EDFA), and frequency-doubled to $780~\mathrm{nm}$ via a Second Harmonic Generation (SHG) crystal (Fig.~\ref{fig:beat_locking}).

Most of this light ($90\%$) is distributed to the MOT optics via a fast-switching AOM. The remaining $10\%$ is mixed with the master laser on a fast ($\sim10\,\mathrm{GHz}$) photodiode to generate beat-note signals for frequency offset locking.

\begin{figure}[h!]
\centering
\includegraphics[width=0.6\textwidth]{Cap_2_figures/beat_peaks.png}
\caption{Example of beat-note signals on a spectrum analyzer, used for frequency offset locking of the cooler and repumper lasers relative to the master laser~\cite{refRazzai}.}
\label{fig:beat_peakes}
\end{figure}

\CCsolved{riferimenti}{The master-cooler and master-repumper beat notes are isolated using RF filters}{La figura con i battimenti sullo spettro non viene mai richiamata nel testo (verificato su tutta la tesi: zero occorrenze della sua \texttt{\textbackslash ref}). Questa frase e' il posto naturale per il rimando, visto che descrive esattamente quello che la figura mostra.}{Fatto: la figura adesso e' richiamata, con \texttt{(Figure \textbackslash ref\{fig:beat\_peakes\})} in fondo a questa frase, e ho verificato che l'etichetta esiste e che il rimando si risolve. Il conto complessivo delle figure mai richiamate scende di uno.} and used as reference signals to stabilize the frequency offsets of the cooler and repumper lasers with respect to the master laser~\cite{refRazzai} (Figure \ref{fig:beat_peakes}).